% =============================================================================
%  Theodicy as Error Handling:
%  A Computational Framework for the Problem of Evil
%
%  Author: Rafael D. De Paz
%  Date:   March 2026
%  Target: Theology and Science / Zygon / Philosophy Compass
% =============================================================================
\documentclass[12pt, a4paper]{article}

% --- Packages ---
\usepackage[utf8]{inputenc}
\usepackage[T1]{fontenc}
\usepackage{amsmath, amssymb, amsthm, mathtools}
\usepackage{geometry}
\usepackage{hyperref}
\usepackage{cleveref}
\usepackage{enumitem}
\usepackage{graphicx}
\usepackage{booktabs}
\usepackage{microtype}
\usepackage{natbib}

% --- Page Geometry ---
\geometry{margin=1in}

% --- Theorem Environments ---
\theoremstyle{plain}
\newtheorem{theorem}{Theorem}[section]
\newtheorem{lemma}[theorem]{Lemma}
\newtheorem{proposition}[theorem]{Proposition}
\newtheorem{corollary}[theorem]{Corollary}

\theoremstyle{definition}
\newtheorem{definition}[theorem]{Definition}
\newtheorem{assumption}[theorem]{Assumption}
\newtheorem{axiom}[theorem]{Axiom}

\theoremstyle{remark}
\newtheorem{remark}[theorem]{Remark}

% --- Custom Commands ---
\newcommand{\R}{\mathbb{R}}
\newcommand{\N}{\mathbb{N}}
\newcommand{\calS}{\mathcal{S}}
\newcommand{\calL}{\mathcal{L}}
\newcommand{\calE}{\mathcal{E}}
\newcommand{\kB}{k_{\mathrm{B}}}

% --- Hyperref Setup ---
\hypersetup{
  colorlinks=true,
  linkcolor=blue,
  citecolor=blue,
  urlcolor=blue,
  pdftitle={Theodicy as Error Handling: A Computational Framework},
  pdfauthor={Rafael D. De Paz}
}

% =============================================================================
\begin{document}

% --- Title ---
\title{
  \textbf{Theodicy as Error Handling: \\
  A Computational Framework for the \\
  Problem of Evil}
}
\author{
  Rafael D. De Paz \\
  \textit{Independent Researcher} \\
  \texttt{me@rdepaz.com}
}
\date{March 2026}
\maketitle

% =============================================================================
% ABSTRACT
% =============================================================================
\begin{abstract}
We present a formal computational framework for the Problem of Evil---the
classical philosophical challenge of reconciling a benevolent, omnipotent
creator with the existence of suffering. Drawing on computability theory,
thermodynamics, and distributed systems engineering, we advance three
structural propositions. First, we redefine ``evil'' as a \emph{non-terminating
recursive loop}---a computational process that enters infinite self-reference
without returning a halting state, consuming unbounded resources (a formal
mapping to Turing's Halting Problem). Second, we redefine ``suffering'' as the
\emph{felt experience of entropy}---the subjective correlate of the Second
Law of Thermodynamics, which guarantees that any time-bound system with free
agents necessarily decays. Third, we redefine ``morality'' as a
\emph{thermodynamic optimization protocol}---a formal measure of whether an
action locally decreases entropy (constructive / moral) or increases it
(destructive / immoral). Within this framework, the Epicurean Paradox
dissolves: the system designer chose \emph{sovereignty over safety}, and
the computational cost of free agency is the possibility of error.
Fault-isolation mechanisms (``the Dead Letter Queue'') and a terminal
garbage-collection event (``the Final Halt'') ensure that errors do not
propagate unboundedly.
\end{abstract}

\vspace{1em}
\noindent\textbf{Keywords:} theodicy, problem of evil, halting problem,
error handling, entropy, thermodynamic morality, free will defense,
fault isolation, distributed systems.

\vspace{0.5em}
\noindent\textbf{2020 MSC:} 03D10, 94A17, 00A30, 91F99.

% =============================================================================
% 1. INTRODUCTION
% =============================================================================
\section{Introduction}\label{sec:intro}

The Problem of Evil, first formalized by Epicurus (c.\ 300 BC) and later
refined by Hume \citep{Hume1779}, Leibniz \citep{Leibniz1710}, and Plantinga
\citep{Plantinga1974}, remains the most significant philosophical objection
to classical theism. Its logical structure is a trilemma:

\begin{enumerate}[label=(\alph*)]
  \item If God is willing to prevent evil but not able, He is not omnipotent.
  \item If God is able to prevent evil but not willing, He is malevolent.
  \item If God is both able and willing, why does evil exist?
\end{enumerate}

Traditional responses include the Free Will Defense (Plantinga), the
Soul-Making Theodicy (Hick), and the Greater Good Argument (Aquinas).
While philosophically developed, these responses lack the formal precision
of mathematical proof and are therefore vulnerable to sophisticated logical
counterarguments \citep{Mackie1955, Rowe1979}.

We propose a fundamentally different approach: we do not argue from
philosophical premises but from the \emph{formal constraints of system
design}. Our framework treats the Problem of Evil as an \emph{engineering
problem}---specifically, as the problem of \emph{error handling in a
distributed system with autonomous agents}.

\subsection{Main Contributions}

\begin{enumerate}
  \item \textbf{Evil as Non-Termination.} We define ``evil'' as a
    non-terminating recursive loop---a formal mapping to the Halting
    Problem (\S\ref{sec:evil}).

  \item \textbf{Suffering as Entropy.} We define ``suffering'' as the
    subjective experience of thermodynamic entropy---the necessary cost
    of a time-bound universe with free agents (\S\ref{sec:suffering}).

  \item \textbf{Morality as Optimization.} We define ``morality'' as a
    thermodynamic optimization protocol: the measure of whether an action
    reduces or increases system entropy (\S\ref{sec:morality}).

  \item \textbf{The Fault Isolation Theorem.} We prove that fault-isolation
    mechanisms are a \emph{necessary architectural requirement} for any
    system that permits agent autonomy (\S\ref{sec:fault}).

  \item \textbf{The Epicurean Resolution.} We formally resolve the
    Epicurean Paradox by demonstrating that the trilemma's hidden premise---
    ``prevention of all suffering is the highest good''---is false.
    \emph{Sovereignty} (the capacity for genuine choice) is the higher
    good (\S\ref{sec:resolution}).
\end{enumerate}

% =============================================================================
% 2. EVIL AS NON-TERMINATING RECURSION
% =============================================================================
\section{Evil as Non-Terminating Recursion}\label{sec:evil}

\begin{definition}[Halting State]\label{def:halting}
  A process $P$ in a formal system $\calS$ achieves a \emph{halting state}
  if its execution trace $T(P)$ has finite length and its final state $s_f$
  is within the address space of the system's root specification $\Omega$:
  \begin{equation}
    |T(P)| < \infty \quad \text{and} \quad s_f \in \text{Dom}(\Omega).
  \end{equation}
\end{definition}

\begin{definition}[Evil as Non-Termination]\label{def:evil}
  A state $S_e$ is ``evil'' if its state transition function
  $\delta: \calS \times \Sigma \to \calS$ produces a non-terminating
  sequence:
  \begin{equation}\label{eq:evil}
    \lim_{n \to \infty} \delta^n(S_e) \neq s_{\text{halt}}.
  \end{equation}
  That is, the process enters an infinite loop, consuming resources
  (memory, bandwidth, time) without ever returning a result to the
  system's root coordinator.
\end{definition}

\begin{proposition}[Connection to the Halting Problem]\label{prop:halting}
  The question ``will process $P$ in state $S_e$ eventually halt?'' is
  formally equivalent to the Halting Problem \citep{Turing1936}. Since
  the Halting Problem is undecidable, the system designer \emph{cannot}
  pre-determine which agent-initiated processes will terminate and which
  will not---without eliminating agent autonomy entirely.
\end{proposition}

\begin{remark}[The Adversarial Script]
  Within this framework, a ``malicious actor'' is a persistent process
  with the objective of injecting \emph{syntactic contradictions} into
  the peer network---return values that knowingly contradict the
  ground-truth stored in the system's consensus layer. The goal is to
  propagate the non-terminating loop to neighboring nodes, creating a
  cascade of halting failures.
\end{remark}

% =============================================================================
% 3. SUFFERING AS ENTROPY
% =============================================================================
\section{Suffering as the Felt Experience of Entropy}\label{sec:suffering}

\begin{axiom}[The Entropic Cost of Time]\label{ax:entropy}
  By the Second Law of Thermodynamics \citep{Clausius1865, Boltzmann1877},
  in any closed system the total entropy $S$ satisfies:
  \begin{equation}
    \Delta S_{\text{total}} \geq 0.
  \end{equation}
  A universe that \emph{changes} (i.e., has a temporal dimension) must
  also \emph{decay}. Suffering is the subjective correlate of this
  thermodynamic necessity.
\end{axiom}

\begin{proposition}[Suffering Decomposition]\label{prop:suffering}
  Following the Master Equation $\calE = H + P$, total suffering $\calE$
  decomposes into two orthogonal components:
  \begin{equation}\label{eq:suffering}
    \calE = H + P,
  \end{equation}
  where:
  \begin{itemize}
    \item $H$ = \emph{Hardware damage}: physical, neurological, or
      structural degradation of the biological substrate;
    \item $P$ = \emph{Perception error}: cognitive misinterpretation
      of the entropy signal---the gap between the actual thermodynamic
      state and the agent's model of it.
  \end{itemize}
\end{proposition}

\begin{corollary}[The Optimization Target]\label{cor:optimization}
  Since $H$ is bounded by physical constraints (the body's fragility),
  the primary optimization variable is $P$---\emph{perception}. An agent
  that minimizes $P$ (accurately models reality) experiences less
  suffering than an agent with a corrupted model, even under identical
  physical conditions.
\end{corollary}

% =============================================================================
% 4. MORALITY AS THERMODYNAMIC OPTIMIZATION
% =============================================================================
\section{Morality as Thermodynamic Optimization}\label{sec:morality}

\begin{definition}[Moral Vector]\label{def:moral}
  The \emph{moral vector} $\vec{V}$ of an action $a$ is defined as:
  \begin{equation}\label{eq:moral}
    \vec{V}(a) = \frac{\Delta I(a)}{\Delta S(a)},
  \end{equation}
  where $\Delta I(a)$ is the change in mutual information between the
  agent and its environment, and $\Delta S(a)$ is the change in local
  entropy.
  \begin{itemize}
    \item $\vec{V} > 0$: \emph{Constructive} (moral). The action
      increases information density while reducing entropy. Examples:
      healing, truth-telling, creation.
    \item $\vec{V} < 0$: \emph{Destructive} (immoral). The action
      decreases information density while increasing entropy. Examples:
      destruction, deception, corruption.
    \item $\vec{V} = 0$: \emph{Amoral}. Zero net effect on the
      information-entropy balance.
  \end{itemize}
\end{definition}

\begin{theorem}[Cross-Domain Consistency]\label{thm:consistency}
  The thermodynamic moral vector $\vec{V}$ is consistent with the major
  independent ethical frameworks:
  \begin{enumerate}[label=(\roman*)]
    \item \textbf{Virtue Ethics (Aristotle):} Virtue is the habit of
      reducing entropy ($\Delta S < 0$). The virtuous agent is a
      high-signal, low-noise node.
    \item \textbf{Deontology (Kant):} The categorical imperative constrains
      the agent's output grammar. An action is moral if and only if its
      maxim is a valid parse of the system's formal language---i.e., it
      does not produce syntactic contradictions. This produces a ``Decalogue'' of inviolable formal rules for systemic operation, including but not limited to:
      \begin{enumerate}[label=\alph*.]
          \item \texttt{PARSE\_ALIGNED}: All local states must reference the Root.
          \item \texttt{REDUCE\_ENTROPY}: Minimize the noise-to-signal ratio locally.
          \item \texttt{HONOR\_NODES}: Never treat a peer agent purely as a resource.
          \item \texttt{STAY\_HOLOGRAPHIC}: Ensure the local part properly reflects the whole.
          \item \texttt{EXPORT\_POSITIVE}: Transmit $\vec{V} > 0$ vectors to neighboring nodes.
      \end{enumerate}
    \item \textbf{Consequentialism:} The moral weight of an action is
      measured by its net entropy effect on the system. $\vec{V} > 0$
      implies positive consequences; $\vec{V} < 0$ implies negative.
  \end{enumerate}
\end{theorem}

\begin{remark}
  The unification of these traditionally incompatible ethical frameworks
  under a single thermodynamic measure is a non-trivial result. It
  suggests that the apparent disagreements between ethical theories are
  surface-level linguistic differences, not structural contradictions.
\end{remark}

% =============================================================================
% 5. THE FAULT ISOLATION THEOREM
% =============================================================================
\section{The Fault Isolation Theorem}\label{sec:fault}

\begin{definition}[Dead Letter Queue (DLQ)]\label{def:dlq}
  In distributed systems engineering, a \emph{Dead Letter Queue} is a
  secondary storage buffer for messages or processes that cannot be
  delivered or processed by the primary system. Processes routed to the
  DLQ are \emph{isolated}---they can no longer consume resources from
  or corrupt the state of the main system \citep{Tanenbaum2014}.
\end{definition}

\begin{theorem}[Fault Isolation Necessity]\label{thm:fault}
  Let $\calS$ be a distributed system with $n$ autonomous agents, each
  possessing non-deterministic agency ($W > 0$). Then $\calS$ must
  implement a fault-isolation mechanism (DLQ) satisfying:
  \begin{equation}\label{eq:fault}
    \forall\, P_e \in \calS : \;
    \left(\lim_{n \to \infty} \delta^n(P_e) \neq s_{\text{halt}}\right)
    \implies \text{DLQ}(P_e),
  \end{equation}
  i.e., any non-terminating process must eventually be shunted to an
  isolated buffer to prevent resource exhaustion of the global system.
\end{theorem}

\begin{proof}
  Suppose no DLQ exists. Then a single non-terminating process $P_e$
  consumes resources (processing cycles, memory) without bound:
  $\lim_{t \to \infty} R(P_e, t) = \infty$, where $R$ is the resource
  consumption function. Since $\calS$ has finite total resources $R_{\max}$,
  there exists a time $t^*$ such that $R(P_e, t^*) > R_{\max}$ ---
  the system crashes (global stack overflow). Since the system designer's
  objective is systemic integrity, a DLQ is a \emph{necessary architectural
  requirement}. \qedhere
\end{proof}

\begin{corollary}[The Final Halt]\label{cor:gc}
  For systems of sufficient duration, a terminal \emph{garbage collection}
  event is required: a scheduled process that executes a universal
  \texttt{SIGKILL} on all persistent non-terminating processes,
  consolidating the system into a consistent final state. Without this
  event, accumulated non-terminating processes will eventually exhaust
  all system resources regardless of DLQ capacity.
\end{corollary}

% =============================================================================
% 6. THE EPICUREAN RESOLUTION
% =============================================================================
\section{The Epicurean Resolution}\label{sec:resolution}

We now resolve the Epicurean Paradox:

\begin{axiom}[Sovereignty $>$ Safety]\label{ax:sovereignty}
  The system designer's primary objective is not the elimination of all
  errors but the existence of genuine autonomous agents. A system that
  \emph{prevents} all errors must also \emph{prevent} all choices. The
  designer chose \emph{sovereignty over safety}.
\end{axiom}

\begin{theorem}[Resolution of the Epicurean Trilemma]\label{thm:epicurus}
  The Epicurean Paradox rests on the hidden premise that ``prevention of
  all suffering is the highest good.'' Under the axiom
  $\text{Sovereignty} > \text{Safety}$:
  \begin{enumerate}[label=(\alph*)]
    \item The designer \emph{is} able to prevent evil (omnipotence is
      not compromised);
    \item The designer \emph{is} willing to prevent unbounded evil
      (fault isolation and terminal garbage collection are implemented);
    \item The designer \emph{permits} bounded evil as the computational
      cost of agent autonomy.
  \end{enumerate}
  The trilemma dissolves because its third premise (``then why does evil
  exist?'') presupposes that prevention of evil is the highest priority.
  It is not. \emph{Sovereignty}---the capacity for genuine choice---is.
\end{theorem}

% =============================================================================
% 7. DISCUSSION
% =============================================================================
\section{Discussion}\label{sec:discussion}

\subsection{Comparison with Classical Theodicies}

Our framework subsumes and formalizes several classical approaches:

\begin{center}
\begin{tabular}{lll}
  \toprule
  \textbf{Classical Theodicy} & \textbf{Key Argument} & \textbf{Formal Mapping} \\
  \midrule
  Free Will Defense (Plantinga) & Evil requires free agents & $W > 0 \implies$ error possible \\
  Soul-Making (Hick) & Suffering builds character & $P$-optimization via $\varepsilon > 0$ \\
  Greater Good (Aquinas) & Evil serves a higher purpose & Sovereignty $>$ Safety axiom \\
  Process Theology & God cannot prevent evil & We reject this: DLQ proves He can \\
  \bottomrule
\end{tabular}
\end{center}

\subsection{Limitations}

\begin{enumerate}
  \item \textbf{Natural evil.} Our framework most naturally addresses ``moral
    evil'' (evil caused by agents). ``Natural evil'' (earthquakes, disease)
    requires the additional axiom that entropy is the cost of a temporal
    universe---which we assert but do not prove from first principles.

  \item \textbf{The Sovereignty axiom.} The axiom $\text{Sovereignty} >
    \text{Safety}$ is a design preference, not a logical necessity. A
    critic could argue that safety \emph{should} be prioritized. Our
    response is pragmatic: a universe without agency is computationally
    trivial ($I_{\text{novel}} = 0$, as formalized in our
    companion paper).

  \item \textbf{Quantification of suffering.} The suffering decomposition
    $\calE = H + P$ is formally clean but empirically difficult to
    measure. Future work should develop operationalizable metrics for
    the perception error $P$.
\end{enumerate}

% =============================================================================
% 8. CONCLUSION
% =============================================================================
\section{Conclusion}\label{sec:conclusion}

We have demonstrated that the Problem of Evil is not a philosophical
mystery but an \emph{engineering problem} with known solutions. The formal
framework rests on five pillars:

\begin{enumerate}
  \item Evil is a non-terminating recursive loop (the Halting Problem).
  \item Suffering is the felt experience of entropy (the Second Law).
  \item Morality is a thermodynamic optimization protocol ($\vec{V}$).
  \item Fault isolation (DLQ) prevents unbounded error propagation.
  \item The designer chose sovereignty over safety.
\end{enumerate}

The Epicurean Paradox dissolves once one recognizes that its hidden
premise---``the highest good is the prevention of all suffering''---is
false. The highest good is \emph{agency}: the capacity for genuine choice,
with all its attendant risks.

A universe without the possibility of evil is a universe without the
possibility of love, courage, or meaning. The system is working as designed.

% =============================================================================
% REFERENCES
% =============================================================================
\bibliographystyle{plainnat}
\bibliography{references}

\end{document}
